\documentclass[UTF8,a4paper,zihao=5]{ctexart}

\usepackage[left=1.5cm,right=1.5cm,top=1.35cm,bottom=1.35cm]{geometry}
\usepackage{graphicx}
\usepackage{booktabs}
\usepackage{tabularx}
\usepackage{array}
\usepackage{float}
\usepackage{subcaption}
\usepackage{enumitem}
\usepackage[hidelinks]{hyperref}
\usepackage{xcolor}
\usepackage{caption}
\usepackage{titlesec}

\setlength{\parindent}{2em}
\setlength{\parskip}{0.08em}
\linespread{0.96}
\setlist[itemize]{leftmargin=2em,itemsep=0.1em,topsep=0.2em}
\setlist[enumerate]{leftmargin=2em,itemsep=0.1em,topsep=0.2em}
\captionsetup{font=small,labelfont=bf,skip=2pt}
\renewcommand{\arraystretch}{1.18}
\titlespacing*{\section}{0pt}{0.55em}{0.25em}
\titlespacing*{\subsection}{0pt}{0.35em}{0.15em}

\newcommand{\placeholder}[2]{%
  \fbox{%
    \begin{minipage}[c][#1][c]{0.94\linewidth}
      \centering\small\color{gray!70!black}#2
    \end{minipage}
  }%
}

\title{\bfseries 星空与萤火：基于 SSVEP 空间注意调制的\\可穿戴脑机接口注意力训练系统}
\author{}
\date{}

\begin{document}
\maketitle
\vspace{-3.2em}

\section{研究背景与问题定位}

注意力障碍是一类跨诊断、跨年龄段的功能问题。ADHD 是最典型的人群之一，世界 ADHD 联盟共识报告指出，ADHD 约影响 5.9\% 的青少年和 2.5\% 的成年人\cite{faraone2021adhd}；美国 CDC 2022 年调查显示，3--17 岁儿童中约 11.4\% 曾被诊断为 ADHD\cite{cdc2024adhd}。除 ADHD 外，孤独症谱系障碍儿童常伴随选择性注意、共同注意和抗干扰困难，世界卫生组织近年估计全球约每 127 人中有 1 人处于孤独症谱系\cite{who2025autism}。卒中后注意功能下降也较常见，注意、执行功能和加工速度障碍会影响患者康复参与、日常活动和重返工作能力\cite{gillespie2015stroke}。对这些人群而言，注意力问题不只是“容易走神”，而是会直接影响学习、康复、社交和生活自理。

现有干预方法各有价值，但仍存在空白。药物治疗对 ADHD 核心症状有效，但并不适合所有用户，且依从性、副作用和长期管理成本需要被认真考虑。认知行为治疗、作业治疗和传统注意训练依赖专业人员持续参与，训练效果也容易受动机影响。EEG 神经反馈提供了一条非药物、可闭环的路线，但多数方案依赖前额叶 beta 功率或 theta/beta 比值作为“专注度”指标。这类指标能够反映一定程度的警觉和认知唤醒，却很难回答更具体的问题：用户的注意力此刻到底指向哪个空间目标。

稳态视觉诱发电位（steady-state visual evoked potential, SSVEP）适合切入这一空白。当用户注视或注意一个以固定频率闪烁的视觉目标时，枕叶视觉皮层会产生同频率及其谐波成分的稳态脑电响应。Morgan 等早在 1996 年即证明，空间选择性注意能够调制 SSVEP 幅值\cite{morgan1996ssvep}；Kim 等进一步指出，注意会带来与同步化相关的 response gain\cite{kim2007gain}。因此，SSVEP 不只是视觉刺激诱发信号，也可作为空间注意分配的读出。不同目标使用不同频率标记后，系统可以估计用户把注意力分给了目标还是干扰项。

本项目选择 SSVEP 作为核心范式，聚焦“空间选择性注意、持续注视、抗干扰和动态追踪”这几类更适合用 SSVEP 量化的能力。项目定位为注意力训练与评估辅助工具，不替代临床诊断和治疗，而是依托赛题一提供的 HNNK 五通道可穿戴脑电头环和 HybridBCI 科研科创平台，构建一个可穿戴、可演示、可逐步验证的脑机接口应用系统。

\section{系统整体设计思路}

本项目拟参加赛题一“可穿戴脑机接口应用系统设计与开发”。系统基于赛方提供的五通道脑电头环与 HybridBCI 平台，调用平台支持的 SSVEP 指令输出，开发一套面向注意力训练的闭环交互应用。系统设计顺序是先确定神经科学问题，再生成交互形式：训练目标不是让用户追逐一个抽象分数，而是让用户反复练习“把注意力稳定投向正确目标，并在干扰出现时保持在那里”。

\begin{figure}[!htbp]
  \centering
  \includegraphics[width=0.62\linewidth]{arch_diagram.png}
  \caption{系统总体架构。系统由脑电采集、SSVEP 指令输出、注意指标计算和场景反馈四部分组成，形成“采集--解码--评估--反馈”的闭环。}
  \label{fig:arch}
\end{figure}

如图 \ref{fig:arch} 所示，头环采集 O1、O2 等枕叶相关通道脑电信号；HybridBCI 平台输出 SSVEP 指令或目标频率识别结果；应用层将结果转换为目标响应强度、目标/干扰竞争指数、持续达标时间和恢复时间；反馈层再把这些指标映射为视觉、声音和任务难度变化。

\subsection{硬件基础}

硬件部分优先使用赛方提供设备，不另行设计脑电采集硬件。HNNK 五通道头环包含 O1、O2、C3、C4 及前额参考/接触电极，其中 O1、O2 位于枕叶区，适合采集视觉皮层相关 SSVEP 信号。设备通过蓝牙 5.0 与上位机连接，官方说明书给出的续航约 8 小时，裸机重量约 85 g，具备可穿戴和现场演示可行性。O1、O2、C3、C4 为氯化银电极，需要配合电极膏并确保穿过头发接触头皮；前额水凝胶电极用于参考与稳定接触。系统将在训练开始前设置标定流程，提示佩戴位置、电极接触和蓝牙连接状态，避免因佩戴不良导致误判。

\begin{figure}[!htbp]
  \centering
  \includegraphics[height=0.19\textheight]{hardware_photo.png}
  \caption{HNNK 五通道可穿戴脑电头环。图中硬件说明支持本项目使用 O1/O2 枕叶通道采集 SSVEP 信号，并体现系统对赛方硬件平台的适配。}
  \label{fig:hardware}
\end{figure}

\subsection{任务设计}

六级任务不是视觉主题的横向扩展，而是认知负荷的纵向递进。第一、二关先让用户在低干扰条件下形成稳定 SSVEP 响应，解决“能否维持”的问题；第三、四关加入空间候选和同场干扰，解决“能否选对并抑制错误目标”的问题；第五、六关再加入运动目标和突发高显著刺激，解决“在复杂环境中能否恢复和保持”的问题。这一路径对应临床训练中常见的梯度：先建立基础持续性注意，再训练选择和转移，最后进入抗干扰与动态控制。

\begin{table}[!htbp]
  \centering
  \caption{六级注意力训练任务与临床/认知功能映射}
  \label{tab:tasks}
  \begin{tabularx}{\linewidth}{>{\raggedright\arraybackslash}p{3.1cm}>{\raggedright\arraybackslash}p{2.1cm}>{\raggedright\arraybackslash}X>{\raggedright\arraybackslash}p{3.1cm}}
    \toprule
    训练阶段 & 界面名称 & 主要训练能力 & SSVEP 设计 \\
    \midrule
    基础持续性注意训练 & 涟漪绽放 & 单目标持续注视、注意保持 & 单一目标频率，低干扰 \\
    视觉耐力与稳定跟随训练 & 萤火引路 & 小目标稳定跟随、视觉疲劳适应 & 多个同频小目标，局部闪烁 \\
    空间搜索与转移性注意训练 & 星图寻迹 & 在多个空间候选中按任务线索转移注意 & 目标频率与背景干扰频率并存 \\
    选择性注意与抑制控制训练 & 真假萤火 & 抑制高显著干扰，维持目标选择 & 目标/干扰竞争指数 \\
    动态目标追踪训练 & 飞燕破云 & 移动目标追踪、错误恢复 & 移动目标加瞬态干扰 \\
    高负荷抗干扰保持训练 & 流星试炼 & 强干扰下的目标保持和冲动抑制 & 高负荷抗干扰任务 \\
    \bottomrule
  \end{tabularx}
\end{table}

训练中的核心指标是目标 SSVEP 响应是否稳定高于干扰响应。对于单目标任务，系统关注目标频率的持续响应和波动；对于双频任务，系统计算竞争指数：目标响应 /（目标响应 + 干扰响应）。当竞争指数稳定较高时，场景变清晰、目标光源更稳定、环境声音更通透；当用户被干扰吸引时，反馈会变弱，但不采用强惩罚。康复训练中，用户更需要可恢复的引导，而不是失败感。

\subsection{交互与界面设计}

界面原则是“训练中少界面，训练前后给信息”。训练前界面用于连接设备、完成标定、选择关卡和训练时长；训练中尽量不出现数值、得分和复杂按钮，只保留必要的暂停/退出；训练后展示平均竞争指数、有效专注时长、分心恢复时间、抗干扰表现和疲劳提示。

本系统不把分数条、能量柱作为主要反馈，而是把 SSVEP 目标嵌入自然场景中的局部发光元素，例如花蕊、萤火、星点、水面高光和飞行光点。这里的美学不是外层包装，而是刺激设计和依从性设计的一部分。若刺激裸露成频闪块，系统更容易调试，却更难长期训练；将刺激安放在自然语义中，反而迫使系统更严格地处理刺激面积、亮度调制、颜色分层和反馈节奏。

\begin{figure}[!htbp]
  \centering
  \begin{subfigure}[t]{0.48\linewidth}
    \centering
    \placeholder{2.15cm}{Landing page 截图占位\\项目定位、设备入口与开始标定按钮}
    \caption{首页用于说明系统定位与赛题匹配。}
  \end{subfigure}
  \hfill
  \begin{subfigure}[t]{0.48\linewidth}
    \centering
    \placeholder{2.15cm}{标定页面截图占位\\O1/O2 通道质量、频率测试与通过提示}
    \caption{标定页用于完成生理接口对齐。}
  \end{subfigure}

  \vspace{0.5em}
  \begin{subfigure}[t]{0.48\linewidth}
    \centering
    \placeholder{2.15cm}{第一关截图占位\\单一 SSVEP 目标、低干扰场景与专注反馈}
    \caption{第一关证明基础持续性注意任务。}
  \end{subfigure}
  \hfill
  \begin{subfigure}[t]{0.48\linewidth}
    \centering
    \placeholder{2.15cm}{第六关截图占位\\目标星/月、流星或极光干扰、无 HUD 反馈}
    \caption{第六关证明高负荷抗干扰任务。}
  \end{subfigure}
  \caption{关键界面设计。四张图分别服务于系统定位、标定可行性、低干扰基线任务和高负荷抗干扰任务的论证。实际提交前可用对应界面截图替换占位框。}
  \label{fig:ui}
\end{figure}

\section{实现路径与可行性}

项目技术路线分为三个阶段（图 \ref{fig:roadmap}）。第一阶段为视觉任务与交互原型验证，已完成基础场景、六级任务和注意力反馈逻辑；第二阶段为赛方设备接入与闭环演示，使用 HNNK 头环完成佩戴和信号质量检查，通过 HybridBCI 平台调用 SSVEP 模块并将输出映射为训练状态；第三阶段为评估和优化，进行健康志愿者小样本测试，验证指令触发稳定性、反馈延迟、视觉疲劳评分和任务难度梯度。

\begin{figure}[!htbp]
  \centering
  \includegraphics[width=0.96\linewidth]{techflow_diagram.png}
  \caption{系统实现路径。方案优先使用赛方平台现有 SSVEP 指令输出，降低接入风险，再逐步扩展目标/干扰竞争指数和训练记录。}
  \label{fig:roadmap}
\end{figure}

阶段性指标包括：设备在常规演示环境下稳定连接；训练反馈延迟尽量控制在 500 ms 内；单次训练时长 5--20 分钟可配置；六个关卡均可演示；训练后记录主观疲劳评分、有效专注时长、分心恢复时间和关卡表现。初期不强行自研完整 SSVEP 解码算法，优先使用赛题允许的现成平台输出，保证方案审查和实物开发的可实现性。

\section{应用价值与软硬件说明}

预期应用场景包括四类：第一，ADHD 和注意困难儿童的辅助训练，重点训练在干扰存在时保持目标注意的能力；第二，卒中和脑损伤后的注意康复，辅助观察患者在左右空间、动态目标和干扰条件下的表现；第三，普通人群的注意健康维护，作为短时专注训练和正念辅助工具；第四，脑机接口教育展示，把脑电采集、SSVEP 识别、实时反馈和应用场景完整呈现给评委和观众。

\begin{table}[!htbp]
  \centering
  \caption{应用系统软硬件说明}
  \label{tab:hardware}
  \begin{tabularx}{\linewidth}{>{\raggedright\arraybackslash}p{2.1cm}>{\raggedright\arraybackslash}p{3.8cm}>{\raggedright\arraybackslash}p{2.2cm}>{\raggedright\arraybackslash}X}
    \toprule
    类型 & 名称 & 来源 & 用途 \\
    \midrule
    EEG 设备 & HNNK 五通道可穿戴脑电头环 & 赛方提供 & 采集 O1/O2 等通道 EEG \\
    BCI 平台 & HybridBCI 科研科创平台 & 赛方提供 & 输出 SSVEP/P300 指令，提供统一通信接口 \\
    终端设备 & 笔记本电脑 & 参赛队自备 & 运行训练系统和展示界面 \\
    显示设备 & 60 Hz 及以上显示器，推荐高刷新率 & 参赛队自备 & 呈现 SSVEP 视觉刺激 \\
    软件系统 & 训练交互应用与本地记录模块 & 自研与开源组件 & 场景渲染、训练逻辑、数据保存与报告展示 \\
    \bottomrule
  \end{tabularx}
\end{table}

\section{总结}

本项目的核心判断是：注意力训练不应只测“专注强不强”，还应尽量测“注意力是否稳定落在正确目标上”。SSVEP 的频率标签机制为这个问题提供了合适入口。依托赛题一的五通道头环和 HybridBCI 平台，本系统可以在少通道、可穿戴条件下完成空间注意训练闭环，并通过场景化反馈改善长期训练的可接受性。

项目边界也需要明确：SSVEP 主要适用于视觉空间注意训练，不能覆盖所有注意功能；系统需要控制闪烁频率、亮度、面积和训练时长，避免视觉疲劳；当前阶段是应用系统设计与原型开发，临床有效性仍需后续伦理审批、小样本试验和对照研究验证。基于这一边界，本项目更适合作为面向注意力训练的可穿戴 BCI 应用原型，为后续实物开发和临床合作留下清晰路径。

\begin{thebibliography}{12}
\bibitem{faraone2021adhd}
Faraone S V, Banaschewski T, Coghill D, et al.
The World Federation of ADHD International Consensus Statement: 208 Evidence-based Conclusions about the Disorder.
\textit{Neuroscience \& Biobehavioral Reviews}, 2021, 128: 789--818.

\bibitem{cdc2024adhd}
Centers for Disease Control and Prevention.
Data and Statistics on ADHD. 2024.
\url{https://www.cdc.gov/adhd/data/index.html}

\bibitem{who2025autism}
World Health Organization.
Autism fact sheet. 2025.
\url{https://www.who.int/news-room/fact-sheets/detail/autism-spectrum-disorders}

\bibitem{gillespie2015stroke}
Gillespie D C, Bowen A, Chung C S, Cockburn J, Knapp P, Pollock A.
Rehabilitation for post-stroke cognitive impairment: an overview of recommendations arising from systematic reviews of current evidence.
\textit{Clinical Rehabilitation}, 2015, 29(2): 120--128.

\bibitem{morgan1996ssvep}
Morgan S T, Hansen J C, Hillyard S A.
Selective attention to stimulus location modulates the steady-state visual evoked potential.
\textit{Proceedings of the National Academy of Sciences}, 1996, 93(10): 4770--4774.

\bibitem{kim2007gain}
Kim Y J, Grabowecky M, Paller K A, Muthu K, Suzuki S.
Attention induces synchronization-based response gain in steady-state visual evoked potentials.
\textit{Nature Neuroscience}, 2007, 10: 117--125.

\bibitem{norcia2015review}
Norcia A M, Appelbaum L G, Ales J M, Cottereau B R, Rossion B.
The steady-state visual evoked potential in vision research: A review.
\textit{Journal of Vision}, 2015, 15(6): 4.

\bibitem{vialatte2010review}
Vialatte F B, Maurice M, Dauwels J, Cichocki A.
Steady-state visually evoked potentials: focus on essential paradigms and future perspectives.
\textit{Progress in Neurobiology}, 2010, 90(4): 418--438.

\bibitem{lin2006cca}
Lin Z, Zhang C, Wu W, Gao X.
Frequency recognition based on canonical correlation analysis for SSVEP-based BCIs.
\textit{IEEE Transactions on Biomedical Engineering}, 2006, 53(12): 2610--2614.

\bibitem{chinchani2022cbmi}
Chinchani A M, Paliwal S, Ganesh S, et al.
Tracking momentary fluctuations in human attention with a cognitive brain-machine interface.
\textit{Communications Biology}, 2022, 5: 1346.

\bibitem{huang2022neurofeedback}
Huang X, Mak J, Wears A, et al.
Using neurofeedback from steady-state visual evoked potentials to target affect-biased attention in augmented reality.
\textit{2022 IEEE Engineering in Medicine and Biology Society}, 2022: 2314--2318.

\bibitem{jeunet2016visual}
Jeunet C, Lotte F, Batail J M, Philip P, Micoulaud-Franchi J A.
Neurofeedback therapy for enhancing visual attention: state-of-the-art and challenges.
\textit{Frontiers in Neuroscience}, 2016, 10: 352.
\end{thebibliography}

\end{document}
